\chapter{Evaluierungskonzept und Ergebnisse}
\label{chap:evaluierungskonzept_und_ergebnisse}

\section{Evaluationsdatensätze}

Die Bewertung der Modelleistung erfolgt getrennt auf Trainings-, Validierungs- und Testdatensätzen.
Diese Aufteilung wird im Rahmen des zuvor beschriebenen Kapitels \ref{chap:datensplit_und_training} vorgenommen.

Der Testdatensatz wird erst nach der finalen Modellauswahl für die abschließenden Evaluierungen verwendet.
Dadurch wird verhindert, dass Informationen aus dem Testdatensatz in die Optimierung des Modells einfließen und die Bewertung verfälschen.

Ein wichtiger Aspekt dabei ist die Behandlung von fachlich Zusammenhängenden Daten. Artikel die denselben Fingerprint besitzen, dürfen nur gemeinsam in einem Teil des Splits vorkommen.
Ein Fingerprint repräsentiert die Struktur und die wesentlichen Merkmale eines Artikels, sodass Artikel mit demselben Fingerprint als fachlich zusammenhängend betrachtet werden.

Dadurch wird verhindert, dass nahezu ähnliche Beispiele gleichzeitig im Trainings- und Testdatensatz vorkommen.

\section{Aufbau der Evaluations-Suite}

Die Evaluation des Learning-To-Rank-Modells erfolgt automatisiert mithilfe der implementierten Evaluations-Suite \cite{src_evalsuite}.

Die Evaluationspipeline besteh dabei aus vier Schirtten:

\begin{table}[htbp] 
  \centering 
  \caption{Aufbau der Evaluations-Suite} 
  \begin{tabular}{p{1cm}p{5cm}p{8cm}} 
    \toprule Nr. & Skript & Zielsetzung \\ 
    \midrule 
    1 & step\_941\_validate\_leakage\_safety.py & Prüfung auf Zielinformations-Leakage \\ 
    2 & step\_942\_validate\_retrieval\_index.py & Bewertung des Retrieval-Index \\ 
    3 & step\_943\_analyze\_candidate\_coverage.py & Analyse der Kandidatenabdeckung \\ 
    4 & step\_944\_analyze\_ltr\_run.py & Analyse des Learning-to-Rank-Modells \\ 
    \bottomrule 
  \end{tabular} 
\end{table}

Die Ergebnisse dieser Pipeline werden strukturiert in den folgenden Verzeichnissen abgelegt:

\begin{itemize}
  \item \texttt{evaluation/step_941_leakage_safety} für die Prüfung auf Zielinformations-Leakage
  \item \texttt{evaluation/step_942_retrieval_validation} für die Bewertung des Retrieval-Index
  \item \texttt{evaluation/step_943_candidate_coverage} für die Analyse der Kandidatenabdeckung
  \item \texttt{evaluation/step_944_ltr_run_analysis} für die Analyse des Learning-to-Rank-Modells
\end{itemize}

\section{Verwendete Rankingmetriken}
 
Bei einem Learning-To-Rank-Modell steht die Qualität der gesamten Rangliste im Vordergrund. 
Für vorliegende Aufgabe ein EPLAN-Fenstermakro\cite{EPLAN-Makros} anhand eines oder mehrer EPLAN-Artikel \cite{EPLANArtikelverwaltung} zu finden,
wurde mehrere komplementäre Rankingmetriken verwendet. 
Hierdurch lassen sich unterschiedliche Aspekte der Modellqualität wie Position des Ziel-Makros \cite{EPLAN-Makros} , Erkennung fachlich verwandter Makros \cite{EPLAN-Makros}, sowie die allgemeine Qualität der Rangfolge beurteilen.

Die Berechnung erfolgt innerhalb des Traingingmoduls und der Evaluationspipeline. \cite{srcStep931,srcStep944}.

\subsection{Target Recall@k} 

Durch den Target Recall@k wird gemessen, ob sich das korrekte Ziel-Makro \cite{EPLAN-Makros} innerhalb der ersten $k$ Elemente einer Rangliste befindet \cite{manning2008ir}.

Dieser berechnet sich wie folgt: 

\begin{align} Recall@k = \frac{ \# \text{Anfragen mit korrektem Ziel in Top-}k } { \# \text{Anfragen} } 
\end{align} 

Die Kennzahl gehört zu den grundlegenden Bewertungsmetriken des Information Retrievals \cite{manning2008ir}.

Der Wertebereich liegt zwischen 0 und 1. Ein Wert von \begin{align} Recall@1 = 0.80 \end{align} bedeutet zum Beispiel, dass bei 80\,\% der Anfragen das korrekte EPLAN-Fenstermakro \cite{EPLAN-Makros} direkt auf Rang 1 liegt. 

Folgenden Varianten werden im Projek verwendet: 

\begin{itemize} 
  \item target\_recall@1 
  \item target\_recall@3 
  \item target\_recall@5 
  \item target\_recall@10 
\end{itemize} 

Diese Metriken sind die wichtigsten Kennzahlen für die Beurteilung der Makrozuordnung.

\subsection{Family Recall@k}

Neben der Erkennung des exakten Ziel-Makros \cite{EPLAN-Makros} wird auch die Erkennung der korrekten Makro-Familie bewertet.
Die Berechnung erfolgt analog zum Target Recall@k und basiert damit auf dem klassischen Recall-Konzept des Information Retrieval \cite{manning2008ir}. Der Unterschied liegt darin, dass jedes EPLAN-Fenstermakro \cite{EPLAN-Makros} als Treffer zählt, das in der gleichen Makro-Familie wie das Ziel-Makro \cite{EPLAN-Makros} vorkommt.
Dadurch kann erkannt werden, ob das System die richtige Makro-Familie erkennt, auch wenn das exakte Ziel-Makro \cite{EPLAN-Makros}  nicht in den Top-k-Ergebnissen enthalten ist.

Hierfür werden folgenden Ausprägungen verwendet:

\begin{itemize} 
  \item family\_recall@1 
  \item family\_recall@3 
  \item family\_recall@5 
  \item family\_recall@10 
\end{itemize}

Die Differenz zwischen dem Family Recall@k und dem Target Recall@k liefert einen Hinweis über die Variantenfehler innerhalb einer Makro-Familie.

\subsection{Mean Reciprocal Rank}

Der Mean Reciprocal Rank (MRR) zählt zu den etablierten Bewertungsmetriken des Information Retrievals und wird bis heute sowohl in klassischen Suchsystemen als auch in modernen Retrieval- und Rankingverfahren eingesetzt \cite{voorhees1999mrr,zhu2024llmirsurvey}.

Dieser berechnet sich wie folgt:

\begin{align} 
  MRR = \frac{1}{N} \sum_{i=1}^{N} \frac{1}{rank_i} 
\end{align} 

mit 

\begin{description} 
  \item[$N$] Anzahl der Anfragen 
  \item[$rank_i$] Rangposition des korrekten Zielmakros für Anfrage $i$ 
\end{description}

Diese Metrik bevorzugt Platzierungen auf den ersten Rängen der Rangliste. Ein EPLAN-Fenstermakro \cite{EPLAN-Makros} auf Rang 1 trägt somit stärker zur MRR bei als ein EPLAN-Fenstermakro \cite{EPLAN-Makros} auf Rang 5.

Dadurch lässt sich die praktische Nutzbarkeit des Ranking-Systens besser bewerten.

\subsection{Hierarchical nDCG}

Für das vorliegende Projekt sind nicht alle Fehler gleich kritisch. Wenn zwei nahezu identische Makros verwechselt werden ist dies weniger schwerwiegend, 
als die Verwechslung von EPLAN-Fenstermakros \cite{EPLAN-Makros} aus unterschiedlichen Familien. Deswegen wird zur Bewertung eine hierachische Bewertungsmetrik verwendet.
Grundlage der Berechnung bildet die von Järvelin und Kekäläinen eingeführte Kennzahl Discounted Cumulative Gain (DCG) \cite{jarvelin2002dcg}.

Die Berechnung ergibt sich wie folgt: 

\begin{align} 
  DCG@k = \sum_{i=1}^{k} \frac{2^{rel_i}-1} {\log_2(i+1)} 
\end{align} 

nDCG@k (normalized Discounted Cumulative Gain) ist die Erweiterung des DCG, die die Bewertung der Rangliste auf einen normierten Bereich [0,1] skaliert.
nDCG zählt heute zu den wichtigsten Bewertungsmetriken für Ranking- und Suchsysteme \cite{jarvelin2002dcg} und wird wie folgt berechnet:

\begin{align} 
  nDCG@k = \frac{DCG@k} {IDCG@k} 
\end{align} 

wobei $IDCG@k$ die ideale Rangliste beschreibt. Diese Metrik wird auch von dem verwendeten Lambda-Rank-Algorithmus genutzt. (siehe Kapitel \ref{chap:datensplit_und_training})


Für die vorliegende Anwendung wird die hierachische Relevanz verwendet: 

\begin{align} 
  rel = \begin{cases} 
    3, & \text{exaktes Zielmakro}\\ 
    2, & \text{Makro derselben Familie}\\ 
    0, & \text{sonst} 
  \end{cases} 
\end{align} 

Dadurch wird ein Ranking vorzugsweise verwendet, welches fachlich näher Alternativen auf den vorderen Positionen platziert.

In diesem Projekt werden die folgenen Metriken verwendet:

\begin{itemize} 
\item hierarchical\_ndcg@1 
\item hierarchical\_ndcg@3 
\item hierarchical\_ndcg@5 
\item hierarchical\_ndcg@10 
\end{itemize}

\subsection{Bedeutung der Metriken innerhalb der Evaluation} 

Die Bewertungsmetriken bilden gemeinsam verschiedene Aspekte der Rankingqualität ab. 

\begin{table}[htbp] 
  \centering 
  \caption{Bedeutung der verwendeten Rankingmetriken} 
  \begin{tabular}{p{4cm}p{9cm}} 
    \toprule Metrik & Bewerteter Qualitätsaspekt \\ \midrule Target Recall@k & Position des exakten Ziel-Makros \cite{EPLAN-Makros} \\ Family Recall@k & Erkennung der korrekten Makrofamilie \\ MRR & Durchschnittlicher Zielrang \\ Hierarchical nDCG & Qualität der gesamten Rangfolge \\ 
    \bottomrule 
  \end{tabular} 
\end{table}


\section{Evaluation der Retrievalphase}

Bevor die Qualtität des Learning-to-Rank-Modells bewertet wird, muss untersucht werden, ob das korrekte Ziel-Makro \cite{EPLAN-Makros} überhaupt durch Retrievalkomponenten gefunden wird.
Ein Rankingmodell kann ein Ziel-Makro \cite{EPLAN-Makros} nur dann korrekt platzieren, wenn es in der Kandidatenmenge vorhanden ist.

Deswegen erfolgt die Evaluation dieser Retrievalphase in einem getrennten Abschnitt.

Diese Analysen werden durch die  Module \texttt{step\_942\_validate\_retrieval\_index.py} und \texttt{step\_943\_analyze\_candidate\_coverage.py} durchgeführt \cite{srcStep942,srcStep943}.

\subsection{Bewertung des Retrieval-Index}

Die Qualität des Retrivalindexes wird durch die erste Evaluationsstufe untersucht. Hierzu werden die Kandidatenlisten die durch das Retrieval erzeugt werden dierkt ausgewertet. 
Dies erfolgt durch \texttt{step\_942\_validate\_retrieval\_index.py} \cite{srcStep942}.

Dabei werden folgende Kennzahlen berechnet:

\begin{itemize} 
\item Target Recall@1 
\item Target Recall@10 
\item Target Recall@100 
\item Family Recall@1 
\item Family Recall@10 
\item Family Recall@100 
\item durchschnittlicher Zielrang 
\item Rangperzentile 
\item Anzahl nicht gefundener Zielmakros 
\end{itemize}

Dadurch kann die grundsätzliche Leistungsfähigkeit der Retrievalkomponenten bewertet werden, um die Eignung der Kandidatenmenge für die nachfolgende Rankingphase zu beurteilen.

Target Recall@100 besitzt hier eine hohe Bedeutung. Wenn ein Ziel-Makro \cite{EPLAN-Makros} nicht unter den Top-100-Kandidaten gefunden wird, kann es durch das spätere Learning-to-Rank-Modell nicht korrekt platziert werden.
Hierdurch wird dann die theoretische maximale Leistung des Learning-to-Rank-Modells festgelegt.

\subsection{Candidate Coverage}

Die Abdeckung der Kandidaten ist neben der Retrievalqualität eine weiter wichtiger Faktor.

Die  Analyse erfolgt dabei über \texttt{step\_943\_analyze\_candidate\_coverage.py} \cite{srcStep943}.

Dabei werden unterschiedliche Retrievalsignale darauf bewertet, wie gut das korrekte Ziel-Makro \cite{EPLAN-Makros} in der Kandidatenmenge aufgenommen wird.

Hierzu werden die einezelen Retrievalsignale getrennt voneinander analysiert:

\begin{itemize} 
  \item BM25 
  \item Exact Match 
  \item Reciprocal Rank Fusion 
  \item kombinierte Retrievalvarianten 
\end{itemize}

Für jede Variante wird dabei bewertet, wie oft das korrekte Ziel-Makro \cite{EPLAN-Makros} in den verschiedenen Kandidatenmengen enthalten ist.

Die Analyse erfolgt dabei für: 

\begin{align} 
  Top_{10} 
\end{align} 

\begin{align} 
  Top_{50} 
\end{align} 

\begin{align} 
  Top_{100} 
\end{align} 

der Kandidaten.

Dadurch kann der Beitrag der einzelnen Retievalkomponenten auf die Kandidatenmenge und der Einfluss der einzelnen Retrievalsignale auf die spätere Rankingqualität bewertet werden.

\subsection{Interpretation der Retrievalergebnisse}

Dies getrennte Evaluation ermöglicht es, Fehlkassifikationen in zwei Kategorien zu unterteilen:

\begin{enumerate} 
  \item Das korrekte Ziel-Makro \cite{EPLAN-Makros} ist nicht Bestandteil der Kandidatenmenge. 
  \item Das korrekte Ziel-Makro \cite{EPLAN-Makros} ist vorhanden, wird jedoch durch das Rankingmodell falsch eingeordnet. 
\end{enumerate}

Im ersten Fall liegt das Problem im Retrievalprozess. Das Rankingmodell kann somit das korrekte Ziel-Makro \cite{EPLAN-Makros} nicht finden.
Im zweiten Fall liegt das Problem im Rankingprozess. Das korrekte Ziel-Makro \cite{EPLAN-Makros} ist zwar in der Kandidatenmenge enthalten, wird jedoch vom Learning-to-Rank-Modell nicht auf die vorderen Ränge eingeordnet.

Dadurch lassen sich die Fehlerquellen im Retrieval- und Rankingprozess voneinander unterscheiden und gezielt verbessern.

\section{Leakage-Sicherheitsprüfung}

Ein zentraler Bestandteil in der Evaluation ist eine Leakage-Sicherheitsprüfung auf Zielinformationen. 
Unter Leakage wird die unbeabsichtigte Nutzung von Informationen verstanden, die zum Zeitpunkt einer realen Vorhersage nicht verfügbar sein dürften und dadurch zu unrealistisch guten Modellergebnissen führen würden \cite{kaufman2012leakage}.
Da dieses System auf Artikel \cite{EPLANArtikelverwaltung} - Makro \cite{EPLAN-Makros} -Zuordnungen trainiert wird, besteht das Risiko, 
das Informationen über das tatsächliche Ziel-Makro \cite{EPLAN-Makros} in die Trainings- und Evaluierungsdaten gelangen können und somit die Aussagekraft der verwendeten Evaluationsmetriken verfälscht wird.

Un dies zu vermeiden, enthält die Pipeline eine Leakage-Prüfung, die vor der eigentlichen Modellbewertung ausgeführt wird .
Die Implementierung erfolgt in \texttt{step\_941\_validate\_leakage\_safety.py} \cite{srcStep941}.

\subsection{Ziele der Leakage-Prüfung}

Die Ziele der Leakage-Prüfung sind:

\begin{itemize} 
\item Identifikation verbotener Zielinformationen innerhalb der Query-Daten 
\item Erkennung direkter Ziel-Makro \cite{EPLAN-Makros} -Referenzen 
\item Prüfung auf Familien- und Zielinformationen innerhalb von Eingangsmerkmalen 
\item Kontrolle der Kandidatenprofile 
\item Überprüfung der Datensplits auf Überschneidungen 
\item Sicherstellung der Trennung von Training, Validierung und Test 
\end{itemize}

Somit wird sichergestellt, dass das Modelltraining nur mit den Informationen erfolgt, die später auch im Produktivbetrieb verfügbar sind.

\subsection{Prüfung der Query-Daten}

Die verwendete Querry-Darstellung repräsentiert einen eingehenden Artikel \cite{EPLANArtikelverwaltung} beziehungsweise die Informationen aus einem Datenblatt, die für eine Suchanfrage an das trainiert Modell wichtig sind.

Wichtig hierbei ist es, dass keine direkten Hinweise auf das Ziel-Makro \cite{EPLAN-Makros} in den Query-Daten enthalten sind, welche unnmittelbar auf dieses verweisen können.

Deswegen untersucht die Leakage-Prüfung folgenden Punkte:

\begin{itemize} 
  \item Zielmakrokennungen 
  \item Zielfamilieninformationen 
  \item direkten Zielmakro-Tokens 
  \item verbotenen Metadatenfeldern 
  \item impliziten Zielhinweisen in den Featurevektoren 
\end{itemize}

Sobald solche Informationen gefunden werden, wird der Lauf als Ungültig markiert.

\subsection{Prüfung der Datensplits}

Neben der Querry-Prüfung, erfolgt auch eine Prüfung des erzeugten Datensplits. Dabei wird überprüft, ob unerlaubte Überschneidungen zwischen Training, Validierung und Test vorhanden sind.

Dabei werden folgende Aspekte überprüft:

\begin{itemize} 
  \item Fingerprint-Gruppen 
  \item Fenstermakrozuordnungen \cite{EPLAN-Makros} 
  \item Kandidatenprofile 
  \item technische Zusammenhangskomponenten 
  \item gruppierte Artikelstrukturen 
\end{itemize}

Das Ziel dieser Prüfung ist es, sicherzustellen, dass keine fachlich identischen Beispiele gleichzeitig in mehreren Datensplits vorhanden sind. Dadurch wird ein künstliche Modell-Verbesserung verhindert,
die die Generalisierungsfähigkeit des Modells beeinträchtigen würde.

\subsection{Prüfung der Kandidatenprofile}

Ein weiter Bereich der Prüfung betrifft die Makro \cite{EPLAN-Makros} - kandidatenprofile.
Da diese Profile aus den Traingsdaten entstehen, muss sichergestellt werden, dass keine Informationen aus den Validierungs- oder Testdaten in die Kandidatenprofile einfließen.

Deswegen werden sämtliche Profilinformationen auf ihre Herkunft geprüft, um sicherzustellen, dass ausschließlich Daten des Trainingssplits zur Erstellung der Kandidatenprofile verwendet werden \cite{srcLeakageGuard}.

\subsection{Evaluierungsergebnis}

Nach der Leakage-Prüfung wird ein struktruierter Sicherheitsbericht erstellt. Der Evaluationslauf gilt nur dann als bestanden, wenn keine Verstöße gegen die definierten Leakage-Regeln festgestellt wurden.
Diese Basis bildet die Voraussetzung für die Interpretierbarkeit der nachfolgenden Rankingmetriken.